\section{Conclusion}\label{sec:conclusion}

The finite-window fold model is now organized by a single structural bridge.
After an explicit affine permutation of residues, each fiber multiplicity is a
Fibonacci-lag difference of the indexed Fibonacci partition function.  This
single-fiber identity transfers the full partition pressure sequence
\((P_q)_{q\ge 0}\) to finite-window Zeckendorf fibers: the collision moments
\(S_q(m)\) grow with exactly the same constants \(\lambda_q\), and the same
constants admit polynomial-size nonnegative matrix realizations.

The main gain of the present version is that these pressures no longer appear
as isolated moment exponents.  Their adjacent differences
\(\Delta_q=P_q-P_{q-1}\) select the canonical thickness bands of the model.
For each fixed \(q\), the \(q\)-canonical law proportional to \(d_m(x)^q\) is
exponentially concentrated on the band \([\Delta_q,\Delta_{q+1}]\), and that
band carries explicit two-sided exponential estimates both for its cardinality
and for its visible mass.  At the zero-temperature end, diverging tilts
concentrate on the single limiting thickness \(\frac12\log\varphi\), which
upgrades the largest-fiber law and the diagonal high-moment asymptotic to a
full canonical concentration theorem.

The visible law under uniform input is exactly the endpoint ensemble
\(\mu_{m,1}\).  Its typical thickness band, its R\'enyi entropy rates, and its
visible-tail estimates are therefore not separate corollaries attached to the
moment theory, but direct projections of the same pressure geometry.  In this
form, finite-window Zeckendorf fibers constitute a discrete thermodynamic
system whose canonical pressures are precisely the Fibonacci partition
constants.  Arithmetic refinements of individual recurrence polynomials and
fixed-resolution inverse moment problems remain interesting, but they are
orthogonal to the structural formalism developed here.
